% !TEX program = lualatex
\documentclass[a4paper,11pt]{ltjsarticle}
\usepackage[margin=25mm,top=24mm,bottom=24mm]{geometry}
\usepackage{longtable,booktabs,array,tabularx}
\usepackage{enumitem}
\usepackage{fancyhdr}
\usepackage{hyperref}
\hypersetup{colorlinks=true,linkcolor=black,urlcolor=black}
\setlength{\parindent}{1em}
\setlength{\parskip}{0.45em}
\pagestyle{fancy}
\fancyhead[L]{第09章 ソフトウェアエンジニアリング管理}
\fancyhead[R]{SWEBOK v4.0a 日本語まとめ}
\fancyfoot[C]{\thepage}
\begin{document}
\begin{center}
{\LARGE\bfseries 第09章 ソフトウェアエンジニアリング管理}\par\vspace{1.2em}
{\large Software Engineering Management の日本語まとめ}\par\vspace{0.8em}
SWEBOK Guide v4.0a Chapter 09 を初学者向けに再構成\par\vspace{0.8em}
作成日：2026年5月12日
\end{center}
\hrule\vspace{1em}
\begin{center}
\fbox{\begin{minipage}{0.94\linewidth}
\textbf{この資料の主張}\par\smallskip
ソフトウェアエンジニアリング管理は、単なる進捗表の更新ではない。要求、日程、費用、品質、リスク、人、測定値、契約、供給者、変更を結び付け、利害関係者に価値あるソフトウェアを届けるための管理活動である。ソフトウェアは物理的な製品と異なり、目に見えず、変更しやすく、運用後も継続的に改善される能力である。したがって、管理者は、計画だけでなく、測定、対話、リスク対応、反復的な学習を通じてプロジェクトを導く必要がある。
\end{minipage}}
\end{center}

\begin{center}
\fbox{\begin{minipage}{0.94\linewidth}
\textbf{この資料の読み方}\par\smallskip
専門用語は原則として日本語で示し、必要に応じて英語を括弧内に併記する。英語名は、元資料、標準、ツール名を確認したいときの手がかりとして残す。初学者が前提知識なしに読めるように、各節では「何のための概念か」「実務では何を見ればよいか」を重視する。文体は、だである調で統一する。
\end{minipage}}
\end{center}

\section{全体概要：この章で学ぶこと}
この章は、ソフトウェアプロジェクトをどう始め、どう計画し、どう実行し、どう監視・制御し、どう終結させるかを扱う。さらに、管理を支える測定プロセスと管理ツールも扱う。

ソフトウェア開発では、顧客の理解が進むにつれて要求が変わる。技術も変わり、チームも人間関係や能力差の影響を受ける。したがって、管理とは「最初に決めた計画を守らせること」だけではない。計画と現実の差を見つけ、根拠に基づいて判断し、必要なら計画や範囲を改訂する活動である。

\begin{center}
\fbox{\begin{minipage}{0.94\linewidth}
\textbf{到達目標}\par\smallskip
\begin{itemize}[leftmargin=2em]
  \item ソフトウェアエンジニアリング管理が、計画、見積り、測定、制御、調整、リスク管理、リーダーシップを含むことを説明できる。
  \item 開始と範囲定義、計画、実行、レビューと評価、終結という管理活動の流れを説明できる。
  \item 予測型と適応型の開発ライフサイクルの違いを、管理上の観点から説明できる。
  \item 工数、日程、費用見積り、資源割当て、リスク管理、品質管理の基本を説明できる。
  \item 測定要求、情報ニーズ、測定量、データ収集、分析、改善の関係を説明できる。
\end{itemize}
\end{minipage}}
\end{center}

学習順序としては、まず「ソフトウェア管理が難しい理由」と「管理と測定の関係」を押さえる。その後、「始める」「計画する」「実行する」「評価する」「終える」「測る」「ツールで支援する」という順に読むとよい。

\begin{center}
\fbox{\begin{minipage}{0.94\linewidth}
\textbf{TODO（図） 第09章の全体地図を入れる。}\par\smallskip
画像生成用プロンプト
白背景、モノクロ、A4本文幅に収まる横長の学習用図解を作成する。中央に「ソフトウェアエンジニアリング管理」を置き、周囲に「開始と範囲定義」「プロジェクト計画」「プロジェクト実行」「レビューと評価」「終結」「測定」「管理ツール」を配置する。矢印は「計画する」「実行する」「測る」「学んで改善する」の循環として表す。日本語ラベル、細線、講義ノート風、余白を広めにする。
\end{minipage}}
\end{center}

\section{最初に必要な用語}
\begin{longtable}{>{\raggedright\arraybackslash}p{0.25\linewidth}>{\raggedright\arraybackslash}p{0.70\linewidth}}
\toprule
用語 & 初学者向けの説明 \\
\midrule
\endfirsthead
\toprule
用語 & 初学者向けの説明 \\
\midrule
\endhead
ソフトウェアエンジニアリング管理 & ソフトウェアプロジェクトを計画し、見積り、測定し、監視し、制御し、調整し、リスクを管理し、チームを導く活動である。英語では Software Engineering Management、略して SEM である。 \\
プロジェクト管理 & 期限、費用、範囲、品質、資源、リスク、関係者などを扱い、プロジェクト目標を達成するための管理である。 \\
測定 & 成果物、プロセス、資源に値やラベルを割り当て、判断に使える情報を作ることである。単に数値を集めるだけではなく、管理判断に結び付ける必要がある。 \\
ソフトウェア開発ライフサイクル & 構想から開発、運用、保守、終了までの流れである。英語では Software Development Life Cycle、略して SDLC である。 \\
予測型 & 要求や設計を早い段階で詳しく定め、詳細な計画に基づいて進めるライフサイクルである。 \\
適応型 & 短い反復を通じて要求や計画を段階的に詳しくし、変化へ対応するライフサイクルである。 \\
成果物 & 仕様書、設計書、検査報告、テスト済みソフトウェアなど、測定可能で確認可能な仕事の産物である。 \\
作業分解構造 & 作業を小さな単位に分解した構造である。英語では Work Breakdown Structure、略して WBS である。 \\
リスク & 目標達成に影響する不確実な事象である。悪い影響だけでなく、機会として良い影響を持つ場合もある。 \\
情報製品 & 測定データを分析し、判断に使える形にしたグラフ、指標、報告、結論である。 \\
\bottomrule
\end{longtable}

\section{導入：ソフトウェア管理が難しい理由}
ソフトウェアエンジニアリング管理は、ソフトウェア製品とソフトウェアエンジニアリングサービスを、効率的かつ効果的に、利害関係者の利益に合うように届けるための作業活動の集合である。活動には、計画、見積り、測定、制御、調整、指導、リスク管理が含まれる。

ソフトウェアプロジェクトは、橋や機械装置のような物理的対象のプロジェクトと似ている部分を持つ。しかし、ソフトウェアには独自性がある。ソフトウェアは無形であり、アルゴリズムやデータ構造として表される。また、変更しやすい。ただし、変更しやすいからといって、望む効果を正しく得ることが簡単という意味ではない。むしろ、変更が容易に見えるため、要求変更や手戻りが頻発しやすい。

\subsection{ソフトウェアプロジェクトを複雑にする典型要因}
\begin{itemize}[leftmargin=2em]
  \item 顧客は、何が必要か、何が実現可能かを最初から明確に知らないことが多い。
  \item 理解が進むほど、新しい要求や変更要求が生じる。
  \item 顧客は、要求変更が設計、実装、テスト、日程、費用に及ぼす影響を過小評価しやすい。
  \item ソフトウェアは可鍛性が高いため、線形工程よりも反復的に構築されることが多い。
  \item ソフトウェアは一度納品して終わりではなく、長期に支援され、継続的に改善される能力である。
  \item ソフトウェア構築では、実装中にも設計判断が行われる。物理製品のように、製造前に設計を完全に固定しにくい。
  \item 創造性と規律の両方が必要であり、この均衡を取ることが難しい。
  \item 技術の変化が速く、開発者の知識更新、ツール選定、保守戦略に影響する。
  \item 長さや重量のような物理的単位をそのまま使えないため、見積り、監視、制御が難しい。
  \item 速度とサイクル時間が重要であり、事業やミッションの要求に合わせて継続的に能力を届ける必要がある。
\end{itemize}
\subsection{三つの管理レベル}
\begin{longtable}{>{\raggedright\arraybackslash}p{0.25\linewidth}>{\raggedright\arraybackslash}p{0.35\linewidth}>{\raggedright\arraybackslash}p{0.35\linewidth}}
\toprule
管理レベル & 何を扱うか & 実務で見ること \\
\midrule
\endfirsthead
\toprule
管理レベル & 何を扱うか & 実務で見ること \\
\midrule
\endhead
組織・基盤管理 & 組織方針、人材育成、標準、ツール基盤、再利用戦略、ポートフォリオを扱う。 & 組織標準、採用・育成方針、共通ツール、過去データ、再利用資産。 \\
プロジェクト管理 & 個別プロジェクトの範囲、日程、費用、品質、資源、リスク、供給者を扱う。 & プロジェクト憲章、計画、WBS、見積り、リスク登録簿、進捗報告。 \\
測定プログラム管理 & 測定要求、データ収集、分析、報告、改善を扱う。 & 測定計画、測定量、ダッシュボード、情報製品、改善記録。 \\
\bottomrule
\end{longtable}

\begin{center}
\fbox{\begin{minipage}{0.94\linewidth}
\textbf{管理と測定の関係}\par\smallskip
測定を伴わない管理は、根拠の弱い判断になりやすい。一方、管理目的のない測定は、数値を集めるだけで終わりやすい。ソフトウェアエンジニアリングでは、管理と測定を組み合わせることで、遅延、費用超過、品質低下、機能不足の兆候を早く見つける。
\end{minipage}}
\end{center}

\begin{center}
\fbox{\begin{minipage}{0.94\linewidth}
\textbf{TODO（図） ソフトウェア管理が難しい理由の因果図を入れる。}\par\smallskip
画像生成用プロンプト
白背景、モノクロ、A4本文幅の因果関係図を作成する。中央に「ソフトウェア管理の難しさ」を置き、周囲から「要求が変わる」「無形で測りにくい」「技術変化が速い」「人間要因が大きい」「実装中にも設計が起こる」「運用後も改善が続く」を矢印で接続する。各要因から「見積り困難」「手戻り」「品質リスク」「日程変動」へ細い矢印を伸ばす。日本語ラベル、講義ノート風。
\end{minipage}}
\end{center}

\section{1 開始と範囲定義}
開始と範囲定義は、ソフトウェアプロジェクトに着手するかどうかを判断し、目的、範囲、実現可能性、権限を確認する活動である。ここで確認すべき中心は、要求が何であり、その要求を満たす取り組みが現実的かどうかである。

\subsection{1.1 要求の決定と交渉}
要求の決定と交渉では、要求獲得、要求分析、要求仕様化、要求妥当性確認を通じて、プロジェクトが何を達成すべきかを合意する。要求は、その後の設計、実装、テスト、受け入れ、保守の基礎になる。

実務では、要求はプロジェクト憲章や高水準の開始文書に記録されることが多い。ここで大切なのは、単に要求を集めることではなく、利害関係者の視点を調整し、優先度、制約、成功条件を明確にすることである。

\begin{center}
\fbox{\begin{minipage}{0.94\linewidth}
\textbf{交渉が必要な理由}\par\smallskip
要求はしばしば衝突する。利用者は機能を増やしたいが、予算や日程には限界がある。セキュリティを高めると利便性や費用に影響することがある。したがって、管理者は、要求、制約、リスクを同じ場に置いて交渉する必要がある。
\end{minipage}}
\end{center}

\subsection{1.2 実現可能性分析}
実現可能性分析は、プロジェクト目的を明確にし、代替案を評価し、技術、資源、費用、倫理、環境、社会技術的な制約のもとで、提案された解決策が妥当かを判断する活動である。

初期のプロジェクト範囲記述、製品範囲、主要成果物、期間制約、必要資源の概算を用意する。ここでいう資源には、人員、予算、設備、開発環境、支援組織、外部供給者、必要能力を持つ人が含まれる。

作業分解構造（WBS）は、仕事を小さな単位へ分解して管理しやすくする文書である。ただし、WBS自体は費用や日程の基準線ではない。費用や日程は、次の計画活動で具体化される。

\begin{longtable}{>{\raggedright\arraybackslash}p{0.25\linewidth}>{\raggedright\arraybackslash}p{0.70\linewidth}}
\toprule
確認項目 & 問い \\
\midrule
\endfirsthead
\toprule
確認項目 & 問い \\
\midrule
\endhead
目的 & このプロジェクトは何を解決するのか。成功を何で判断するのか。 \\
範囲 & 何を作り、何を作らないのか。既存システムや外部組織との境界はどこか。 \\
技術 & 必要な技術は利用可能か。実績があるか。未知の技術リスクは何か。 \\
資源 & 必要な人員、スキル、設備、予算、期間は確保できるか。 \\
代替案 & 自作、購入、外注、SaaS利用、既存資産再利用のどれが適切か。 \\
倫理・環境・社会 & 利用者、社会、環境、法規制に悪影響はないか。 \\
\bottomrule
\end{longtable}

\begin{center}
\fbox{\begin{minipage}{0.94\linewidth}
\textbf{TODO（図） 実現可能性分析の観点図を入れる。}\par\smallskip
画像生成用プロンプト
白背景、モノクロ、A4本文幅の放射状図を作成する。中央に「実現可能性分析」を置き、周囲に「目的」「範囲」「技術」「資源」「費用」「倫理」「環境」「社会技術的制約」「代替案」を配置する。各要素に短い説明を添える。日本語ラベル、講義ノート風。
\end{minipage}}
\end{center}

\subsection{1.3 要求のレビューと改訂のプロセス}
要求と範囲は、開発中に変化する。したがって、開始時点で「どのように見直し、どのように変更を受け入れるか」を合意しておく必要がある。

変更は、無秩序に受け入れてはならない。変更管理、影響分析、トレードオフ、反復のふりかえり、マイルストーンレビューなどの手順を定める。変更を受け入れる場合には、要求、設計、コード、テスト、文書への影響を追跡可能性で確認する。さらに、リスク分析を行い、費用、日程、品質、安全性、セキュリティへの影響を見積もる。

\begin{center}
\fbox{\begin{minipage}{0.94\linewidth}
\textbf{範囲は石に刻まれたものではない}\par\smallskip
範囲と要求は、最初に決めたら絶対に変えてはならないものではない。しかし、変更の影響を見ずに変えると、プロジェクトは制御不能になる。重要なのは「変えないこと」ではなく「管理された形で変えること」である。
\end{minipage}}
\end{center}

\section{2 ソフトウェアプロジェクト計画}
ソフトウェアプロジェクト計画は、プロジェクトを成功させるために、ライフサイクル、成果物、見積り、資源、リスク、品質、計画管理を決める活動である。

計画で最初に重要なのは、適切なソフトウェア開発ライフサイクルを選び、必要に応じて調整することである。ライフサイクルは、予測型から適応型までの連続体として考えるとよい。

\begin{center}
\fbox{\begin{minipage}{0.94\linewidth}
\textbf{TODO（図） 予測型から適応型までのライフサイクル連続体を入れる。}\par\smallskip
画像生成用プロンプト
白背景、モノクロ、A4本文幅の横長図を作成する。横軸の左端を「予測型」、右端を「適応型」とする。左側に「詳細要求」「詳細計画」「マイルストーン中心」、右側に「漸進的要求」「短い反復」「継続的な利害関係者関与」を置く。中央に「増分型」「スパイラル型」を配置する。下に「プロジェクトに応じて選び調整する」と注記する。日本語ラベル、講義ノート風。
\end{minipage}}
\end{center}

\subsection{2.1 プロセス計画}
プロセス計画では、どのライフサイクルと方法を使うかを選ぶ。予測型ライフサイクルでは、詳細な要求、詳細な計画、比較的少ない反復を前提にする。適応型ライフサイクルでは、要求の出現や変更を前提にし、短い開発サイクルで計画を更新する。

代表的なライフサイクルには、ウォーターフォール型、増分型、スパイラル型、各種アジャイル型がある。近年は、開発、セキュリティ、運用を統合する Dev/Sec/Ops も重要である。Dev/Sec/Ops は、計画、開発、ビルド、テスト、リリース、配備、運用、監視の各段階で、自動化、継続的監視、セキュリティ適用を行う文化と方法である。

\begin{longtable}{>{\raggedright\arraybackslash}p{0.25\linewidth}>{\raggedright\arraybackslash}p{0.35\linewidth}>{\raggedright\arraybackslash}p{0.35\linewidth}}
\toprule
観点 & 予測型 & 適応型 \\
\midrule
\endfirsthead
\toprule
観点 & 予測型 & 適応型 \\
\midrule
\endhead
要求 & 早い段階で詳細化する。 & 反復を通じて漸進的に詳細化する。 \\
計画 & 初期に詳細計画を作る。 & 初期計画を置き、実行しながら更新する。 \\
リスク低減 & 詳細計画と既知の制約で減らす。 & 短い反復、早期フィードバック、段階的な学習で減らす。 \\
利害関係者関与 & マイルストーンごとの関与が多い。 & 継続的な関与が多い。 \\
向く状況 & 要求が安定し、規制や契約が強い場合。 & 要求の不確実性が高く、迅速な学習が必要な場合。 \\
\bottomrule
\end{longtable}

プロセス計画では、プロジェクトで用いる方法、ツール、標準も選ぶ。要求管理、設計、実装、テスト、構成管理、品質管理、保守、進捗管理などに使うツールは、技術面だけでなく管理面からも検討する必要がある。

\subsection{2.2 成果物の決定}
成果物とは、プロジェクト活動から生まれる測定可能で確認可能な仕事の産物である。ソフトウェアアーキテクチャ記述、設計文書、検査報告、テスト済みソフトウェア、運用手順、移行計画などが該当する。

成果物を決めるときは、過去プロジェクトの部品再利用、市販製ソフトウェア、オープンソースソフトウェア、外部委託、SaaS の利用可能性も評価する。成果物を自作するのか、取得するのか、供給者に作らせるのかによって、計画、リスク、契約、品質確認の方法が変わる。

\subsection{2.3 工数・日程・費用見積り}
見積りは、ソフトウェアプロジェクトで誤りやすい活動である。理由は、ソフトウェア開発の工数が、人の経験、能力、チーム相互作用、組織文化、要求変化、技術変化に大きく左右されるからである。

一つの方法だけで見積もるのは危険である。過去データに基づく推定モデル、類似プロジェクトとの比較、作業者によるボトムアップ見積り、専門家判断などを組み合わせ、差を調整する必要がある。

\begin{longtable}{>{\raggedright\arraybackslash}p{0.25\linewidth}>{\raggedright\arraybackslash}p{0.70\linewidth}}
\toprule
見積り対象 & 実務で見ること \\
\midrule
\endfirsthead
\toprule
見積り対象 & 実務で見ること \\
\midrule
\endhead
工数 & どの作業に、どのスキルの人が、どれだけ必要かを見る。 \\
日程 & タスク依存関係、並行可能性、開始・終了時期、反復回数を考える。 \\
費用 & 人件費、ツール、設備、ライセンス、取得費、外注費、運用準備費を含める。 \\
不確実性 & 見積りを一点ではなく範囲で考え、前提条件とリスクを明記する。 \\
合意形成 & 利害関係者と、利用可能な資源と期限について合意する。 \\
\bottomrule
\end{longtable}

\begin{center}
\fbox{\begin{minipage}{0.94\linewidth}
\textbf{見積りは約束ではなく仮説である}\par\smallskip
初期見積りは、限られた情報から作る仮説である。要求が変わり、技術リスクが明らかになれば、見積りも更新される。重要なのは、前提条件、根拠、不確実性を明示し、測定データで補正していくことである。
\end{minipage}}
\end{center}

\subsection{2.4 資源割当て}
資源割当ては、人、予算、設備、時間、ツール、外部サービスなどを、プロジェクト活動に割り当てる活動である。ソフトウェアでは、人数だけでなく、必要な技能、経験、役割、チーム構成、コミュニケーション経路を考える必要がある。

人を増やせば常に早く終わるわけではない。新しい人を受け入れるための教育、調整、レビューが増えることがある。したがって、資源割当てでは、単純な人数ではなく、チームが実際に価値を生み出せる能力を評価する。

\subsection{2.5 リスク管理}
リスク管理は、目標達成に影響する不確実性を見つけ、分析し、対応し、監視する活動である。ソフトウェアでは、要求変化、技術の未知性、人員不足、供給者遅延、セキュリティ脆弱性、品質不足、統合失敗などが典型的なリスクである。

\begin{longtable}{>{\raggedright\arraybackslash}p{0.25\linewidth}>{\raggedright\arraybackslash}p{0.70\linewidth}}
\toprule
リスク管理活動 & 内容 \\
\midrule
\endfirsthead
\toprule
リスク管理活動 & 内容 \\
\midrule
\endhead
識別 & 起こり得る問題や機会を列挙する。 \\
分析 & 発生確率、影響度、検出しやすさ、優先度を評価する。 \\
対応計画 & 回避、低減、転嫁、受容、予備計画を決める。 \\
監視 & リスクの兆候、しきい値、責任者、対応状況を追跡する。 \\
学習 & 実際に起きたことを記録し、次の見積りやプロセス改善に反映する。 \\
\bottomrule
\end{longtable}

\subsection{2.6 品質管理}
品質管理は、品質を目標として扱い、その目標を満たすための手順、責任、評価方法を計画する活動である。品質は、最後のテストでだけ作られるものではない。要求、設計、実装、レビュー、構成管理、測定、改善の全体を通じて作り込む。

品質要求には、信頼性、可用性、性能、使用性、保守性、セキュリティ、安全性などがある。受け入れ可能なしきい値を定め、技術レビュー、検査、デモ、テスト、監査の方法を計画する。Dev/Sec/Ops の文脈では、セキュリティや品質確認を早い段階へ移す、いわゆるシフトレフトが重要である。

\subsection{2.7 計画管理}
計画管理は、プロジェクト計画と関連する支援計画を、監視し、レビューし、報告し、必要に応じて改訂する活動である。古い予測型プロジェクトでは正式な計画文書が重視される。一方、適応型プロジェクトでは、詳細文書よりも、制御に必要な情報、戦略、方針、手順、追跡可能性を保持することが重視される。

重要なのは、文書の量ではなく、計画が管理に役立つことである。計画は、現実に対して検査され、測定データによって補正され、利害関係者に共有される必要がある。

\section{3 ソフトウェアプロジェクト実行}
プロジェクト実行では、計画を実施し、計画に含まれるプロセスを動かす。選択したライフサイクルに従いつつ、利害関係者の要求を満たし、プロジェクト目標を達成することを狙う。実行の中心には、監視、制御、報告がある。

\subsection{3.1 計画の実施}
プロジェクト活動は、プロジェクト計画と支援計画に従って行われる。活動は、人員、技術、資金などの資源を使い、設計、コード、テストケース、報告書などの作業成果物を生成する。

実務では、計画どおりに進めるだけでなく、計画と現実の差を観察し、必要な修正を行う。実行中に現れる問題は、要求の不足、見積り誤差、技術課題、チーム課題、外部依存の遅れなど多様である。

\subsection{3.2 ソフトウェア取得と供給者契約管理}
ソフトウェア取得とは、自作以外の方法でソフトウェアやサービスを得ることである。取得形態には、市販製ソフトウェア、オープンソース、委託開発、顧客提供ソフトウェア、SaaS、データセット取得などがある。

\begin{longtable}{>{\raggedright\arraybackslash}p{0.25\linewidth}>{\raggedright\arraybackslash}p{0.70\linewidth}}
\toprule
取得形態 & 管理上の注意点 \\
\midrule
\endfirsthead
\toprule
取得形態 & 管理上の注意点 \\
\midrule
\endhead
市販製ソフトウェア & ライセンス、バージョン、保守期限、サポート範囲、脆弱性対応を確認する。 \\
オープンソース & ライセンス適合、依存関係、コミュニティ活性度、セキュリティ更新を確認する。 \\
委託開発 & 作業範囲、成果物、品質基準、納期、検収条件、知的財産、変更手順を契約で明確にする。 \\
SaaS & 可用性、データ保護、サービスレベル、終了時のデータ移行、障害時対応を確認する。 \\
データセット取得 & データ利用権、品質、更新頻度、個人情報、統合方法を管理する。 \\
\bottomrule
\end{longtable}

どの取得形態でも、各構成要素が要求やアーキテクチャと整合するか、統合後に正しく機能するか、意図された運用環境で目的を満たすかを確認する必要がある。供給者契約では、作業範囲、成果物、納期、品質要求、遅延時の扱い、知的財産権、データの扱いを明確にする。

\begin{center}
\fbox{\begin{minipage}{0.94\linewidth}
\textbf{TODO（図） ソフトウェア取得の種類と契約管理の図を入れる。}\par\smallskip
画像生成用プロンプト
白背景、モノクロ、A4本文幅の分類図を作成する。中央に「ソフトウェア取得」を置き、周囲に「市販製ソフトウェア」「委託開発」「オープンソース」「顧客提供」「SaaS」「データセット取得」を配置する。各分類から「ライセンス」「品質要求」「セキュリティ」「統合」「知的財産」「納期」へ点線でつなぐ。日本語、講義ノート風。
\end{minipage}}
\end{center}

\subsection{3.3 測定プロセスの実施}
計画された測定プロセスは、プロジェクト実行中に実施する。目的は、関連性があり、役に立つデータを集めることである。測定は、後述する「測定プロセス計画」と「測定プロセス実行」に基づいて行う。

\subsection{3.4 監視プロセス}
監視プロセスでは、範囲、計画、支援計画への適合を継続的に確認する。各タスクの成果物と完了基準も評価する。工数消費、日程遵守、費用、資源使用状況、品質測定値、リスク状態を分析する。

差異分析は、計画値と実績値のずれを調べる活動である。費用超過、日程遅延、欠陥密度の悪化、脆弱性の増加などを早期に発見する。しきい値を超えたときには、関係者へ報告し、制御プロセスにつなげる。

\subsection{3.5 制御プロセス}
制御プロセスでは、監視結果に基づいて判断し、必要な対応を行う。対応には、是正処置、追加活動、計画改訂、要求仕様改訂、プロトタイピングの導入、追加テスト、リスク対応、場合によってはプロジェクト中止が含まれる。

変更を行うときは、ソフトウェア構成管理と構成制御の手順に従う。判断を記録し、関係者に伝え、計画を必要に応じて改訂し、関連データを残す。

\begin{center}
\fbox{\begin{minipage}{0.94\linewidth}
\textbf{TODO（図） 監視・制御・報告のフィードバックループを入れる。}\par\smallskip
画像生成用プロンプト
白背景、モノクロ、A4本文幅の循環図を作成する。「実行」から「測定データ収集」「監視」「差異分析」「制御判断」「計画改訂」「実行」へ戻る矢印を描く。右側に「報告」を置き、監視と制御判断から利害関係者へ矢印を伸ばす。しきい値超過を警告アイコンで示す。日本語、講義ノート風。
\end{minipage}}
\end{center}

\subsection{3.6 報告}
報告は、合意された時点で、組織内と外部利害関係者に向けて行う。報告は、詳細な内部ステータスをそのまま流すのではなく、対象者の情報ニーズに合わせる必要がある。

\begin{longtable}{>{\raggedright\arraybackslash}p{0.25\linewidth}>{\raggedright\arraybackslash}p{0.70\linewidth}}
\toprule
報告先 & 知りたいこと \\
\midrule
\endfirsthead
\toprule
報告先 & 知りたいこと \\
\midrule
\endhead
プロジェクトチーム & 当面の作業、課題、依存関係、阻害要因、次の反復での対応。 \\
管理委員会 & 日程、費用、リスク、範囲変更、意思決定が必要な事項。 \\
顧客・利用者 & 要求充足状況、デモ可能な成果、受け入れ判断に必要な情報。 \\
供給者 & 契約上の成果物、品質基準、変更要求、連携タスク。 \\
\bottomrule
\end{longtable}

\section{4 ソフトウェアレビューと評価}
レビューと評価は、定められた時点または必要な時点で、プロジェクトが目的へ進んでいるか、利害関係者の要求を満たしているか、プロセス、人員、ツール、方法が有効かを確認する活動である。

\subsection{4.1 要求充足の判定}
利害関係者の満足を達成することは、ソフトウェアエンジニアリング管理者の主要な目標である。進捗は、主要マイルストーン、ソフトウェア技術レビュー、反復サイクルの完了、製品増分の完成時に評価する。

要求からの差異を見つけ、適切な対応を取る。対応には、要求見直し、設計変更、追加テスト、受け入れ基準の明確化、範囲調整が含まれる。構成管理手順に従い、判断と変更の理由を記録する。

\subsection{4.2 パフォーマンスのレビューと評価}
プロジェクト人員の定期的なパフォーマンスレビューは、計画やプロセスに従える可能性、チーム内の困難、対立、支援の必要性を把握するために役立つ。

プロジェクトで用いる方法、ツール、技法も、有効性と適切性の観点で評価する。プロセスは、状況に対して意味があるか、役に立つか、効果があるかを定期的に確認し、必要に応じて変更する。

\section{5 終結}
終結は、プロジェクト全体、主要フェーズ、または反復サイクルが完了し、完了基準が満たされたことを確認する活動である。終結後には、アーカイブ、ふりかえり、プロセス改善を行う。

\subsection{5.1 終結判定}
終結は、指定されたタスクが完了し、完了基準の満足が確認されたときに起こる。ソフトウェア要求が満たされたかどうか、目標達成の程度はどの程度かを判定する。

終結プロセスには、関連する利害関係者を含める。受け入れを文書化し、既知の問題を記録する。未解決の問題を隠して終結すると、運用・保守で重大な混乱を生む。

\subsection{5.2 終結活動}
終結が確認されたら、プロジェクト資料を、合意された方法、保管場所、保管期間に従って保存する。機密情報、ソフトウェア、記録媒体の破棄が必要な場合もある。

組織の測定データベースには、関連するプロジェクトデータを追加する。さらに、プロジェクト、フェーズ、反復のふりかえりを行い、遭遇した問題、リスク、機会を分析する。教訓は、組織の学習と改善に接続する。

\begin{center}
\fbox{\begin{minipage}{0.94\linewidth}
\textbf{終結は「終わった」ことを証明する活動である}\par\smallskip
終結とは、単に開発チームが手を離すことではない。要求充足、成果物、受け入れ、残課題、記録、データ、教訓を確認し、次の運用・保守・改善へ引き渡す活動である。
\end{minipage}}
\end{center}

\section{6 ソフトウェアエンジニアリング測定}
ソフトウェアエンジニアリング測定は、よい管理とよいエンジニアリング実践の基礎である。測定は、組織、プロジェクト、プロセス、作業成果物に適用できる。この節では、プロジェクト、プロセス、作業成果物の水準で測定を扱う。

測定プロセスは、測定要求を定め、測定を計画し、データを収集・分析し、情報製品として伝達し、測定そのものを評価して改善する流れである。

\begin{center}
\fbox{\begin{minipage}{0.94\linewidth}
\textbf{TODO（図） 情報ニーズから情報製品までの測定モデルを入れる。}\par\smallskip
画像生成用プロンプト
白背景、モノクロ、A4本文幅の流れ図を作成する。左から「組織目標」「測定要求」「情報ニーズ」「測定量」「データ収集」「分析」「情報製品」「意思決定」を箱で並べ、矢印で接続する。下段に「測定プロセス評価」と「改善」を置き、情報製品と測定要求へ戻る矢印を描く。日本語ラベル、講義ノート風。
\end{minipage}}
\end{center}

\subsection{6.1 測定コミットメントの確立と維持}
測定を成功させるには、組織とプロジェクトが測定にコミットする必要がある。測定要求は、組織目標とプロジェクト目標から導く。たとえば「市場投入を早める」という目標があれば、サイクル時間、リードタイム、リリース頻度などの測定が必要になる。

測定の範囲も定める。対象は、単一プロジェクト、単一部門、拠点、組織全体などである。時間的範囲も重要である。見積りモデルを較正するためには、過去から現在までの時系列データが必要になることがある。

\begin{longtable}{>{\raggedright\arraybackslash}p{0.25\linewidth}>{\raggedright\arraybackslash}p{0.70\linewidth}}
\toprule
コミットメント項目 & 説明 \\
\midrule
\endfirsthead
\toprule
コミットメント項目 & 説明 \\
\midrule
\endhead
測定要求の確立 & 測定は組織目標とプロジェクト目標に基づく。 \\
測定範囲の確立 & どの組織単位、期間、プロジェクト、成果物を測るか決める。 \\
チームの合意 & 測定の目的、使い方、責任を明確にし、形式的に支援する。 \\
資源の割当て & 担当者、分析者、データ管理者、ツール、教育、予算を確保する。 \\
\bottomrule
\end{longtable}

\subsection{6.2 測定プロセスの計画}
測定プロセスの計画では、情報ニーズを明確にする。情報ニーズとは、利害関係者が判断するために必要な情報である。たとえば、納期が守れるか、品質が十分か、リスクが増えていないか、要求が満たされているかなどである。

次に、情報ニーズから測定量を選ぶ。測定量には、欠陥数、欠陥密度、要求変更数、未解決リスク数、消費工数、完了率、リードタイム、ビルド失敗率、テスト合格率などがある。測定量は、収集可能で、解釈可能で、行動につながるものでなければならない。

\begin{longtable}{>{\raggedright\arraybackslash}p{0.25\linewidth}>{\raggedright\arraybackslash}p{0.35\linewidth}>{\raggedright\arraybackslash}p{0.35\linewidth}}
\toprule
情報ニーズ & 測定量の例 & 使い方 \\
\midrule
\endfirsthead
\toprule
情報ニーズ & 測定量の例 & 使い方 \\
\midrule
\endhead
日程は守れそうか & 完了タスク数、残作業量、リードタイム、速度。 & 計画の更新、資源調整、範囲調整に使う。 \\
品質は十分か & 欠陥密度、再発欠陥数、テスト合格率、失敗率。 & 追加レビュー、追加テスト、リリース判断に使う。 \\
リスクは増えていないか & 未解決リスク数、高優先度リスク数、発生兆候。 & リスク対応、予備計画、意思決定に使う。 \\
要求は満たされているか & 要求充足率、未実装要求数、変更要求数。 & 受け入れ判断、範囲調整、優先順位見直しに使う。 \\
\bottomrule
\end{longtable}

\subsection{6.3 測定プロセスの実行}
測定プロセスの実行では、データを収集し、検証し、保存し、分析する。データ収集は、自動化できる場合がある。たとえば、課題管理、構成管理、テスト管理、CI/CD、監視ツールからデータを取得できる。

データは、目的に応じた厳密さで集計、変換、記録する。分析の結果は、グラフ、数値、指標、結論、推奨事項として情報製品になる。情報製品は、利害関係者が行動できる形で伝える必要がある。

\begin{center}
\fbox{\begin{minipage}{0.94\linewidth}
\textbf{測定値を罰に使わない}\par\smallskip
測定は、プロジェクトを理解し、改善するための道具である。人を罰する目的で使うと、データ隠し、過剰な最適化、指標の形だけの達成を招く。測定の目的、利用者、解釈方法を共有することが重要である。
\end{minipage}}
\end{center}

\subsection{6.4 測定の評価}
測定の評価では、情報製品と測定プロセスが、指定された評価基準を満たしているかを確認する。測定利用者からのフィードバック、内部プロセス、外部監査を使うことがある。

改善候補には、指標の形式変更、測定単位の変更、分類の見直し、収集方法の自動化、データ品質の改善などがある。改善候補の費用と便益を判断し、測定プロセスの所有者と利害関係者へ提案する。

\section{7 ソフトウェアエンジニアリング管理ツール}
管理ツールは、ソフトウェアエンジニアリング管理プロセスの可視性と制御を支援する。近年は、プロジェクトを通じて、計画、収集、記録、監視、制御、報告を統合するツール群が使われることが多い。

\begin{longtable}{>{\raggedright\arraybackslash}p{0.25\linewidth}>{\raggedright\arraybackslash}p{0.70\linewidth}}
\toprule
ツール分類 & 主な用途 \\
\midrule
\endfirsthead
\toprule
ツール分類 & 主な用途 \\
\midrule
\endhead
計画・追跡ツール & 工数・費用見積り、日程作成、WBS、ガントチャート、マイルストーン、反復、アクション項目を追跡する。 \\
リスク管理ツール & リスク識別、分析、監視、確率分布、意思決定木、モンテカルロシミュレーションを支援する。 \\
コミュニケーションツール & 通知、議事録、変更承認、合意形成、進捗共有、関係者への一貫した情報提供を支援する。 \\
測定ツール & 構成管理、テスト、課題、ビルド、運用監視からデータを集め、報告やダッシュボードを作る。 \\
\bottomrule
\end{longtable}

ツールは目的ではなく手段である。ツールを導入しても、測定要求、情報ニーズ、責任、解釈方法が不明確であれば、管理は改善しない。まず「誰が、何を判断するために、どの情報を必要とするか」を決める必要がある。

\begin{center}
\fbox{\begin{minipage}{0.94\linewidth}
\textbf{TODO（図） 管理ツールの情報連携図を入れる。}\par\smallskip
画像生成用プロンプト
白背景、モノクロ、A4本文幅のシステム概念図を作成する。左に「課題管理」「構成管理」「テスト管理」「CI/CD」「運用監視」を縦に配置し、中央の「測定・分析基盤」へ矢印でつなぐ。右に「計画」「リスク管理」「報告」「意思決定」を配置する。日本語ラベル、講義ノート風。
\end{minipage}}
\end{center}

\section{8 トピックと参考文献の対応}
この表は、SWEBOK が示す「トピック対参考資料の行列」を、学習用に簡略化したものである。詳しく学ぶときは、各トピックに対応する文献を読む。

\begin{longtable}{>{\raggedright\arraybackslash}p{0.25\linewidth}>{\raggedright\arraybackslash}p{0.70\linewidth}}
\toprule
トピック & 主な参考資料 \\
\midrule
\endfirsthead
\toprule
トピック & 主な参考資料 \\
\midrule
\endhead
1 開始と範囲定義 & Fairley, Managing and Leading Software Projects; Sommerville, Software Engineering \\
1.1 要求の決定と交渉 & Fairley; Software Requirements KA \\
1.2 実現可能性分析 & Sommerville \\
1.3 要求のレビューと改訂 & Fairley; Software Configuration Management KA \\
2 プロジェクト計画 & Fairley; Boehm and Turner; Sommerville \\
2.1 プロセス計画 & Fairley; Boehm and Turner \\
2.2 成果物の決定 & Fairley \\
2.3 工数・日程・費用見積り & Fairley; Software Engineering Economics KA \\
2.4 資源割当て & Fairley \\
2.5 リスク管理 & Fairley; Boehm and Turner \\
2.6 品質管理 & Fairley; Sommerville; Software Quality KA \\
2.7 計画管理 & Fairley \\
3 プロジェクト実行 & Fairley; Sommerville \\
3.2 取得と供給者契約管理 & Fairley \\
3.4 監視プロセス & Fairley \\
3.5 制御プロセス & Fairley \\
4 レビューと評価 & Fairley; Sommerville \\
5 終結 & PMBOK Guide; Software Extension to PMBOK Guide \\
6 測定 & Practical Software Measurement; ISO/IEC/IEEE 15939 \\
7 管理ツール & Fairley; Practical Software Measurement \\
\bottomrule
\end{longtable}

\section{9 発展的な読み物}
\begin{longtable}{>{\raggedright\arraybackslash}p{0.25\linewidth}>{\raggedright\arraybackslash}p{0.70\linewidth}}
\toprule
文献 & 読むと得られること \\
\midrule
\endfirsthead
\toprule
文献 & 読むと得られること \\
\midrule
\endhead
PMBOK Guide & プロジェクト統合、範囲、日程、費用、品質、資源、コミュニケーション、リスク、調達、利害関係者管理を体系的に学べる。 \\
Software Extension to the PMBOK Guide & 一般的なプロジェクト管理をソフトウェアプロジェクトへ適用するときの考慮点を学べる。 \\
Fairley, Managing and Leading Software Projects & ソフトウェアプロジェクトの計画、監視、制御、リーダーシップを実務的に学べる。 \\
Sommerville, Software Engineering & ソフトウェア工学全体の基礎と、プロジェクト管理・品質・プロセスとの関係を学べる。 \\
Boehm and Turner, Balancing Agility and Discipline & アジャイル性と規律のバランスを、プロジェクト特性に応じて考える方法を学べる。 \\
Practical Software Measurement & 測定要求、情報ニーズ、測定量、情報製品を結び付ける測定プロセスを学べる。 \\
Defense Innovation Board, Software Is Never Done & ソフトウェアを継続的に改善される能力として捉える考え方を学べる。 \\
IEEE 2675-2021 DevOps & 信頼性とセキュリティを含む DevOps の標準的な考え方を学べる。 \\
\bottomrule
\end{longtable}

\section{10 参考文献}
\begin{enumerate}[leftmargin=2em]
  \item Project Management Institute, A Guide to the Project Management Body of Knowledge, 5th ed., 2013.
  \item Software Extension to the Project Management Body of Knowledge, Fifth Edition, Project Management Institute, 2013.
  \item R. E. Fairley, Managing and Leading Software Projects, Wiley IEEE Computer Society Press, 2009.
  \item I. Sommerville, Software Engineering, 10th ed., Addison-Wesley, 2016.
  \item B. Boehm and R. Turner, Balancing Agility and Discipline: A Guide for the Perplexed, Addison-Wesley, 2003.
  \item IEEE, IEEE Standard Adoption of ISO/IEC 15939:2007 Systems and Software Engineering Measurement Process, 2017.
  \item J. McGarry et al., Practical Software Measurement: Objective Information for Decision Makers, Addison-Wesley Professional, 2001.
  \item J. McDonald, Managing the Development of Software-Intensive Systems, Wiley, 2010.
  \item Practical Software and Systems Measurement Continuous Iterative Development Measurement Framework Parts 1-3, Version 2.1, 2021.
  \item S. Sheard et al., “Book Club” Guides a Working Group to Create INCOSE System-Software Interface Products, INSIGHT, 2021.
  \item K. Nidiffer, C. Woody, and T. A. Chick, Program Manager’s Guidebook for Software Assurance, CMU/SEI-2018-SR-025, 2018.
  \item R. E. Fairley, Systems Engineering of Software-Enabled Systems, Wiley, 2019.
  \item Defense Innovation Board, Software Is Never Done: Refactoring the Acquisition Code for Competitive Advantage, 2019.
  \item IEEE Standard 2675-2021, DevOps: Building Reliable and Secure Systems Including Application Build, Package, and Deployment, 2021.
  \item M. Chemuturi and T. Cagley, Mastering Software Project Management: Best Practices, Tools and Techniques, J. Ross Publishing, 2010.
\end{enumerate}
\section{11 画像 TODO 一覧}
本文に配置した画像 TODO を再掲する。実際に図を作る場合は、各 TODO のプロンプトをそのまま画像生成に使うと、A4 資料に合う図を作りやすい。

\begin{enumerate}[leftmargin=2em]
  \item 第09章の全体地図。
  \item ソフトウェア管理が難しい理由の因果図。
  \item 実現可能性分析の観点図。
  \item 予測型から適応型までのライフサイクル連続体。
  \item ソフトウェア取得の種類と契約管理の図。
  \item 監視・制御・報告のフィードバックループ。
  \item 情報ニーズから情報製品までの測定モデル。
  \item 管理ツールの情報連携図。
\end{enumerate}
\section{12 章末確認問題}
\begin{enumerate}[leftmargin=2em]
  \item ソフトウェアエンジニアリング管理は、単なる進捗管理ではない。理由を説明せよ。
  \item ソフトウェアが物理製品と異なるために、管理が難しくなる点を三つ挙げよ。
  \item 管理と測定を切り離すと、どのような問題が起きるかを説明せよ。
  \item 開始と範囲定義で、要求の決定と交渉が重要な理由を説明せよ。
  \item 実現可能性分析で確認すべき観点を五つ挙げよ。
  \item 予測型と適応型のライフサイクルの違いを、要求、計画、利害関係者関与の観点で説明せよ。
  \item 工数・日程・費用見積りがソフトウェアで難しい理由を説明せよ。
  \item リスク管理で行う識別、分析、対応、監視の違いを説明せよ。
  \item 品質管理をテストだけに任せてはいけない理由を説明せよ。
  \item ソフトウェア取得で、契約に明記すべき事項を挙げよ。
  \item 監視プロセスと制御プロセスの違いを説明せよ。
  \item 終結活動で、教訓を組織の学習へつなげることが重要な理由を説明せよ。
  \item 情報ニーズ、測定量、情報製品の関係を説明せよ。
\end{enumerate}
\begin{center}
\fbox{\begin{minipage}{0.94\linewidth}
\textbf{解答の方向性}\par\smallskip
1 は、計画、見積り、測定、制御、調整、リスク管理、リーダーシップを含む点に注目する。2 は、要求変更、無形性、可鍛性、技術変化、人間要因、継続的改善を挙げる。3 は、根拠のない判断と目的のない数値集めを対比する。4 以降は、本文の「合意」「実現可能性」「範囲」「測定」「監視と制御」「学習と改善」という語を使って説明できる。
\end{minipage}}
\end{center}

\end{document}